\section{Qubits, phase, and the Bloch sphere}

\subsection{A qubit is not just an unknown bit}

A classical bit is exactly 0 or 1. A qubit may be in any normalized
superposition $\alpha\ket0+\beta\ket1$. This does not mean that the qubit
secretly chose a classical value before measurement. Relative amplitudes can
interfere, producing effects no ordinary random bit can reproduce.

\subsection{Global phase and relative phase}

Multiplying an entire state by a unit complex number does not change physical
predictions:
\[
\ket{\psi}\quad\text{and}\quad e^{i\gamma}\ket{\psi}
\]
have the same measurement probabilities and are considered equivalent up to a
\emph{global phase}. Relative phase is different. The states
\[
\ket+=\frac{\ket0+\ket1}{\sqrt2},\qquad
\ket-=\frac{\ket0-\ket1}{\sqrt2}
\]
give equal 0/1 probabilities if measured immediately, yet H sends
$\ket+\mapsto\ket0$ and $\ket-\mapsto\ket1$. Their relative sign is observable
through interference.

\subsection{Bloch-sphere coordinates}

After removing an unobservable global phase, every pure one-qubit state can be
written with two real angles:
\[
\ket{\psi}=
\cos\left(\frac{\theta}{2}\right)\ket0+
e^{i\phi}\sin\left(\frac{\theta}{2}\right)\ket1,
\]
where $0\leq\theta\leq\pi$ and $0\leq\phi<2\pi$. These angles define a point on
the unit Bloch sphere:
\[
x=\sin\theta\cos\phi,\qquad
y=\sin\theta\sin\phi,\qquad
z=\cos\theta.
\]

\begin{center}
\begin{tikzpicture}[scale=2.05]
  \shade[ball color=QpuBlue!10,opacity=.32] (0,0) circle (1);
  \draw[QpuLine] (0,0) ellipse (1 and .32);
  \draw[QpuLine,dashed] (0,.98) arc (90:270:.32 and .98);
  \draw[QpuLine] (0,-.98) arc (-90:90:.32 and .98);
  \draw[-{Latex[length=2mm]},thick] (0,0)--(0,1.28) node[above] {$z,\ \ket0$};
  \draw[-{Latex[length=2mm]},thick] (0,0)--(0,-1.28) node[below] {$-z,\ \ket1$};
  \draw[-{Latex[length=2mm]},thick] (0,0)--(1.25,-.25) node[right] {$x,\ \ket+$};
  \draw[-{Latex[length=2mm]},thick] (0,0)--(-1.05,-.45) node[left] {$y,\ \ket{+i}$};
  \draw[-{Latex[length=2.5mm]},very thick,QpuRed] (0,0)--(.62,.72)
    node[above right] {$\ket{\psi}$};
  \draw[QpuRed] (0,.42) arc (90:49:.42) node[midway,above] {\small$\theta$};
  \fill[QpuRed] (.62,.72) circle (.035);
\end{tikzpicture}
\end{center}

The north pole is $\ket0$ and the south pole is $\ket1$. The equator contains
equal-magnitude superpositions. $\ket+$ points along $+x$,
$\ket-=(\ket0-\ket1)/\sqrt2$ along $-x$,
and $\ket{+i}=(\ket0+i\ket1)/\sqrt2$ along $+y$.

\subsection{How gates move a Bloch vector}

Ignoring global phase, many one-qubit gates are rotations:
\begin{itemize}
  \item X rotates the sphere by $\pi$ around the x-axis, exchanging north and
  south.
  \item Y rotates by $\pi$ around the y-axis.
  \item Z rotates by $\pi$ around the z-axis, changing relative phase.
  \item S and T rotate around z by $\pi/2$ and $\pi/4$.
  \item PHASE rotates around z by its supplied angle $\theta$.
  \item H rotates the coordinate description so the z basis and x basis are
  exchanged: $\ket0\leftrightarrow\ket+$ and
  $\ket1\leftrightarrow\ket-$.
\end{itemize}
The sphere is a teaching model for one pure qubit. It cannot display the full
state of entangled multi-qubit systems as one ordinary three-dimensional arrow.

